\documentclass[11pt,a4paper]{article}

% Paquetes esenciales
\usepackage{cmap} % texto copiable/legible por ATS; debe ir antes de fontenc
\usepackage[utf8]{inputenc}
\usepackage[T1]{fontenc}
\usepackage{lmodern}
\usepackage[spanish]{babel}
\usepackage[left=1.5cm, right=1.5cm, top=1.5cm, bottom=1.5cm]{geometry}
\usepackage{hyperref}
\usepackage{enumitem}
\usepackage{titlesec}
\usepackage{xcolor}

% Configuración de enlaces
\hypersetup{
    colorlinks=true,
    linkcolor=blue,
    filecolor=magenta,
    urlcolor=blue,
    pdftitle={CV - Breyner Ismael Ciro Otero},
}

% Formato de secciones (Línea debajo del título)
\titleformat{\section}
  {\large\bfseries\uppercase}
  {}
  {0em}
  {}[\titlerule]
\titlespacing*{\section}{0pt}{1.5ex plus 1ex minus .2ex}{1ex plus .2ex}

% Desactivar numeración de páginas
\pagestyle{empty}

\begin{document}

% --- ENCABEZADO ---
\begin{center}
    {\huge \textbf{BREYNER ISMAEL CIRO OTERO}} \\ \vspace{0.15cm}
    {\large Ingeniero de Software FullStack $\cdot$ TypeScript $\cdot$ Vue/React $\cdot$ Node.js $\cdot$ GCP/AWS $\cdot$ IA aplicada} \\ \vspace{0.15cm}
    Bogotá, Colombia | +57 3183420438 | \href{mailto:breynerismael@gmail.com}{breynerismael@gmail.com} \\ \vspace{0.1cm}
    \href{https://www.linkedin.com/in/breyner-ciro-38011330b/}{LinkedIn} | \href{https://github.com/breynerciro}{GitHub} | \href{https://breynerciro.github.io}{Portafolio}
\end{center}

\vspace{0.2cm}

% --- PERFIL PROFESIONAL ---
\section*{Perfil Profesional}
Ingeniero de Sistemas y Computación (Universidad Nacional de Colombia) que construye producto SaaS de punta a punta: frontend en Vue y React con TypeScript, backends en Node.js y Python, y datos en Firestore, PostgreSQL y BigQuery sobre GCP y AWS. Hoy integra LLMs (Claude, Gemini) en flujos reales de negocio en una plataforma logística. Trabaja con pruebas automatizadas y despliegues en contenedores.

% --- HABILIDADES TÉCNICAS ---
\section*{Habilidades Técnicas}
\begin{itemize}[leftmargin=*, noitemsep]
    \item \textbf{Lenguajes:} TypeScript, JavaScript (ES6+), Python, SQL.
    \item \textbf{Frontend \& Mobile:} Vue 3 (Pinia), React, React Native, Expo, Tailwind CSS, Vite.
    \item \textbf{Backend \& APIs:} Node.js, Express, FastAPI, Django, APIs RESTful, Webhooks, Gmail API, Microsoft Graph.
    \item \textbf{IA aplicada:} Integración de LLMs (Claude vía Vertex AI, Gemini), extracción de datos de documentos, asistentes con datos en vivo.
    \item \textbf{Datos \& Cloud:} Firebase (Firestore, Cloud Functions), BigQuery, Cloud Run, PostgreSQL, MySQL, MongoDB, Redis, AWS, Docker, Nginx, GitHub Actions.
    \item \textbf{Prácticas:} Testing (Vitest, JMeter), Code Reviews, GitFlow, Scrum, Arquitectura de Microservicios.
\end{itemize}

% --- EXPERIENCIA PROFESIONAL ---
\section*{Experiencia Profesional}

\noindent
\textbf{Kila} | \textit{Desarrollador Junior de Ingeniería de Software} \hfill Octubre, 2026 – Actualidad \\
\textit{SaaS logístico para importadores y freight forwarders.}
\begin{itemize}[leftmargin=*, noitemsep, topsep=2pt]
    \item Formo parte del equipo de ingeniería de una plataforma SaaS que centraliza la operación de comercio exterior de sus clientes; aporto de punta a punta en frontend, backend serverless y datos (\textbf{Vue 3, TypeScript, Firebase, BigQuery, Cloud Run}).
    \item Llevo la IA a funciones de producto: asistentes sobre LLMs (\textbf{Claude en Vertex AI, Gemini}) conectados a datos reales del negocio, diseñados para responder solo con información verificable.
    \item Fortalezco la confiabilidad y la oportunidad de los datos operativos y aduaneros con los que los clientes toman decisiones; reduje a la mitad el tiempo de sincronización con el data warehouse (de más de 2 min a cerca de 1 min).
    \item Conecto la plataforma con el ecosistema de correo de los clientes (\textbf{Gmail, Microsoft Graph}) y sostengo la calidad con pruebas automatizadas en 4 servicios.
\end{itemize}

\vspace{0.2cm}

\noindent
\textbf{CardioDiagnostico del Sur} | \textit{Desarrollador Frontend} \hfill Enero, 2024 – Agosto, 2024
\begin{itemize}[leftmargin=*, noitemsep, topsep=2pt]
    \item Construí el sitio institucional de una IPS de cardiología como Single Page Application con \textbf{React} y \textbf{TypeScript} (servicios, estados financieros y contacto), con código fuertemente tipado.
    \item Maqueté una interfaz 100\% responsiva con \textbf{Tailwind CSS} y gestioné el paso a producción con \textbf{Vite} y despliegue en Hostinger.
    \item Facilité el agendamiento de citas con un botón de WhatsApp (mensaje prellenado) y un formulario de contacto con backend propio en \textbf{Node.js/Express} que envía los mensajes por correo.
\end{itemize}

\vspace{0.2cm}

\noindent
\textbf{Proyectos Corporativos (Freelance / U. Nacional)} | \textit{Desarrollador de Software} \hfill 2022 – 2024
\begin{itemize}[leftmargin=*, noitemsep, topsep=2pt]
    \item Integré SPAs con backends relacionales y validé su comportamiento bajo concurrencia con pruebas de carga y estrés.
    \item Estandaricé los entornos de desarrollo y producción con configuraciones reproducibles en \textbf{Docker} y docker-compose.
    \item Automaticé los flujos de integración y entrega continua con canalizaciones en \textbf{GitHub Actions}.
\end{itemize}

% --- PROYECTOS DESTACADOS ---
\section*{Proyectos Destacados}

\noindent
\textbf{HEMODENT – Aplicación Móvil Educativa en Salud (Proyecto de Tesis)} \hfill 2026
\begin{itemize}[leftmargin=*, noitemsep, topsep=2pt]
    \item Desarrollé una app multiplataforma (iOS/Android) para pacientes con hemofilia con \textbf{React Native}, \textbf{TypeScript} y \textbf{Expo} (SDK 57, Expo Router), con persistencia offline (AsyncStorage), animaciones (Reanimated) y recordatorios mediante notificaciones locales.
    \item Diseñé una interfaz (UX/UI) adaptativa por rango de edad, con acceso a protocolos de emergencia dental y un centro de conocimiento médico especializado.
\end{itemize}

\vspace{0.2cm}

\noindent
\textbf{WORKIO – Sistema de Reservas de Coworking (Full Stack)} \hfill 2025
\begin{itemize}[leftmargin=*, noitemsep, topsep=2pt]
    \item Construí con un equipo Scrum un sistema web de reservas de espacios de coworking: API REST en \textbf{Node.js}, \textbf{Express} y \textbf{TypeScript} organizada en capas, y frontend en \textbf{React}, \textbf{Vite} y \textbf{Tailwind CSS} publicado en Netlify.
    \item Implementé el esquema relacional en \textbf{PostgreSQL}, autenticación con \textbf{JWT} y \textbf{bcrypt}, control de acceso por roles y validación de entradas con express-validator.
\end{itemize}

\vspace{0.2cm}

\noindent
\textbf{OSP – Overwatch Sentinel Platform (Plataforma Edge de Visión Artificial, proyecto en equipo)} \hfill 2025
\begin{itemize}[leftmargin=*, noitemsep, topsep=2pt]
    \item Desarrollé una plataforma de vigilancia con visión por computadora usando \textbf{TensorFlow Lite} en \textbf{Raspberry Pi} (Edge Computing) para detección de objetos y optimización de grabación de video.
    \item Diseñé una arquitectura de microservicios con backend (\textbf{Python/Node.js}), frontend (\textbf{React}) y bases de datos (\textbf{PostgreSQL, MongoDB}), comunicados mediante RESTful APIs, OAuth2 (Google) y JWT.
    \item Orquesté el despliegue en 5 entornos con \textbf{Docker}, un API Gateway (\textbf{Vert.x}) con patrón Gatekeeper y proxy reverso (\textbf{Nginx}) con red privada segmentada contra DDoS, SQLi y MitM; validé alta disponibilidad y tolerancia a fallos con pruebas de estrés en \textbf{JMeter}.
\end{itemize}

% --- EDUCACIÓN ---
\section*{Educación}

\noindent
\textbf{Ingeniero de Sistemas y Computación} | Universidad Nacional de Colombia \hfill Graduado: Septiembre 2026 \\
\textit{Bogotá D.C., Colombia.}

\vspace{0.2cm}

\noindent
\textbf{Técnico Laboral por Competencias en Sistemas Informáticos} | POLINTEC \hfill Diciembre 2016 \\
\textit{Puerto Asís, Putumayo.}

% --- CERTIFICACIONES ---
\section*{Certificaciones}
\begin{itemize}[leftmargin=*, noitemsep]
    \item \textbf{Diplomado en Arquitectura y Calidad de Software (192h, 3 módulos):} Universidad Nacional de Colombia. Terminados: Ingeniería de Software Basada en la Nube y Calidad de Software. En curso: Arquitectura de Software a Gran Escala (finaliza en noviembre de 2026).
    \item \textbf{AWS Academy Graduate - Cloud Foundations (20h):} Credly (Noviembre 2025).
    \item \textbf{Curso Básico de Power BI (8h):} Datly Educación (Julio 2025).
\end{itemize}

% --- IDIOMAS ---
\section*{Idiomas}
\noindent
\textbf{Inglés:} B1 | \textbf{Portugués:} B1

\end{document}
